\documentclass[../../main.tex]{subfiles}

\begin{document}

\section{一致可积性: Vitali 收敛定理}\label{section:ch04-一致可积}

通过对有限测度集上可积的函数建立积分号下取极限的合理性准则, 我们结束关于 Lebesgue 积分的第一章. 该准则由以下的引理和命题给出.

\begin{lemma}\label{lemma:ch04-有限集族}
令 $E$ 为有限测度集而 $\delta>0$. 则 $E$ 是有限集族的不交并, 该族中的每个集合的测度小于 $\delta$.
\end{lemma}

\begin{proof}
根据测度的连续性,
\begin{align*}
\lim_{n\to\infty}m(E\sim[-n,n])=m(\varnothing)=0.
\end{align*}
选取自然数 $n_{0}$, 使得 $m(E\sim[-n_{0},n_{0}])<\delta$. 通过选取 $[-n_{0},n_{0}]$ 的足够好的分划, 将 $E\cap[-n_{0},n_{0}]$ 表示为有限集族的不交并, 每个集合的测度小于 $\delta$.
\end{proof}

\begin{proposition}\label{proposition:ch04-绝对连续积分}
令 $f$ 为 $E$ 上的可测函数. 若 $f$ 在 $E$ 上可积, 则对每个 $\varepsilon>0$, 存在 $\delta>0$ 使得
\begin{align}
\text{若 }A\subseteq E\text{ 是可测的且 }m(A)<\delta,\text{则}\int_{A}|f|<\varepsilon.\label{align:ch04-绝对连续}
\end{align}
反过来, 在 $m(E)<\infty$ 的情形, 若对每个 $\varepsilon>0$, 存在 $\delta>0$ 使\eqref{align:ch04-绝对连续}成立, 则 $f$ 在 $E$ 上可积.
\end{proposition}

\begin{proof}
分别对 $f$ 的正部和负部证明定理. 我们因此假定在 $E$ 上 $f\geqslant 0$. 首先假设 $f$ 在 $E$ 上可积. 令 $\varepsilon>0$. 根据积分在区域上的可加性与 Chebyshev 不等式, 若 $A\subseteq E$ 是可测的而 $c>0$, 则
\begin{align*}
\int_{A}f=\int_{\{x\in A\mid f(x)<c\}}f+\int_{\{x\in A\mid f(x)\geqslant c\}}f\leqslant c\cdot m(A)+\frac{1}{c}\int_{E}f.
\end{align*}
选取 $c>0$ 使得 $1/c\cdot\int_{E}f<\varepsilon/2$. 则
\begin{align*}
\int_{A}f<c\cdot m(A)+\varepsilon/2.
\end{align*}
定义 $\delta=\varepsilon/(2c)$. 则对 $\delta$ 的这个选择\eqref{align:ch04-绝对连续}成立.

反过来, 假定 $m(E)<\infty$, 且对每个 $\varepsilon>0$, 存在 $\delta>0$ 使得\eqref{align:ch04-绝对连续}成立. 令 $\delta_{0}>0$ 对应于 $\varepsilon=1$ 的挑战. 由于 $m(E)<\infty$, 根据前一个引理, 我们可以将 $E$ 表示为有限个不交可测子集族 $\{E_{k}\}_{k=1}^{N}$ 的并集, 每个集合的测度小于 $\delta_{0}$. 因此
\begin{align*}
\sum_{k=1}^{N}\int_{E_{k}}f<N.
\end{align*}
根据积分在区域上的可加性, 得到若 $h$ 为具有有限支撑的非负可测函数且在 $E$ 上 $0\leqslant h\leqslant f$, 则 $\int_{E}h<N$. 因此 $f$ 是可积的.
\end{proof}

\begin{definition}[一致可积]\label{definition:ch04-一致可积}
$E$ 上的可测函数族 $\mathcal{F}$ 称为在 $E$ 上\textbf{一致可积的}, 若对每个 $\varepsilon>0$, 存在 $\delta>0$ 使得对每个 $f\in\mathcal{F}$,
\begin{align}
\text{若 }A\subseteq E\text{ 是可测的且 }m(A)<\delta,\text{则}\int_{A}|f|<\varepsilon.\label{align:ch04-一致可积定义}
\end{align}
\end{definition}

\begin{example}\label{example:ch04-一致可积例}
令 $g$ 为 $E$ 上的非负可积函数. 定义
\begin{align*}
\mathcal{F}=\{f\mid f\text{ 在 }E\text{ 上可测且在 }E\text{ 上 }|f|\leqslant g\}.
\end{align*}
则 $\mathcal{F}$ 是一致可积的. 这可以从以下方面得到: 由命题\cref{proposition:ch04-绝对连续积分} ($f$ 换作 $g$), 以及观察到对 $E$ 的任何可测子集 $A$, 根据积分的单调性, 若 $f$ 属于 $\mathcal{F}$, 则
\begin{align*}
\int_{A}|f|\leqslant\int_{A}g.
\end{align*}
\end{example}

\begin{proposition}\label{proposition:ch04-有限族一致}
令 $\{f_{k}\}_{k=1}^{n}$ 为有限函数族, 该族中的每个函数在 $E$ 上可积. 则 $\{f_{k}\}_{k=1}^{n}$ 是一致可积的.
\end{proposition}

\begin{proof}
令 $\varepsilon>0$. 对 $1\leqslant k\leqslant n$, 根据命题\cref{proposition:ch04-绝对连续积分}, 存在 $\delta_{k}>0$ 使得
\begin{align}
\text{若 }A\subseteq E\text{ 是可测的且 }m(A)<\delta_{k},\text{则}\int_{A}|f_{k}|<\varepsilon.\label{align:ch04-单函数估计}
\end{align}
定义 $\delta=\min\{\delta_{1},\cdots,\delta_{n}\}$. 该 $\delta$ 对应于关于族 $\{f_{k}\}_{k=1}^{n}$ 一致可积准则中的 $\varepsilon$ 挑战.
\end{proof}

\begin{proposition}\label{proposition:ch04-极限可积}
假设 $E$ 有有限测度. 令 $\{f_{n}\}$ 为 $E$ 上的一致可积函数序列. 若在 $E$ 上 a.e. 逐点 $\{f_{n}\}\to f$, 则 $f$ 在 $E$ 上可积.
\end{proposition}

\begin{proof}
令 $\delta_{0}>0$ 对应于关于序列 $\{f_{n}\}$ 一致可积准则的 $\varepsilon=1$ 挑战. 由于 $m(E)<\infty$, 根据引理\cref{lemma:ch04-有限集族}, 我们可以将 $E$ 表示为可测子集的有限族 $\{E_{k}\}_{k=1}^{N}$ 的不交并集, 使得对 $1\leqslant k\leqslant N$, $m(E_{k})<\delta_{0}$. 对任何 $n$, 根据积分的单调性和在区域上的可加性,
\begin{align*}
\int_{E}|f_{n}|=\sum_{k=1}^{N}\int_{E_{k}}|f_{n}|<N.
\end{align*}
我们从 Fatou 引理推出
\begin{align*}
\int_{E}|f|\leqslant\liminf\int_{E}|f_{n}|\leqslant N.
\end{align*}
因此 $|f|$ 在 $E$ 上可积.
\end{proof}

\begin{theorem}[Vitali 收敛定理]\label{theorem:ch04-Vitali}
令 $E$ 为有限测度集. 假定函数序列 $\{f_{n}\}$ 在 $E$ 上一致可积.
\begin{align*}
\text{若在 }E\text{ 上 a.e. 逐点 }\{f_{n}\}\to f,\text{则 }f\text{ 在 }E\text{ 上可积且}\lim_{n\to\infty}\int_{E}f_{n}=\int_{E}f.
\end{align*}
\end{theorem}

\begin{proof}
命题\cref{proposition:ch04-极限可积}告诉我们 $f$ 在 $E$ 上可积, 因此, 根据命题\cref{proposition:ch04-可积a.e.有限}, 它在 $E$ 上 a.e. 有限. 因此, 又一次用命题\cref{proposition:ch04-可积a.e.有限}, 通过可能地从 $E$ 删除一个测度为零的集合, 我们假定收敛在整个 $E$ 上是逐点的, 而 $f$ 是实值的. 我们从积分的比较判别法与积分的线性、单调性以及在区域上的可加性推出, 对 $E$ 的任何可测子集 $A$ 和任何自然数 $n$,
\begin{align*}
\left|\int_{E}f_{n}-\int_{E}f\right|=\left|\int_{E}(f_{n}-f)\right|\leqslant\int_{E}|f_{n}-f|=\int_{E\sim A}|f_{n}-f|+\int_{A}|f_{n}-f|\leqslant\int_{E\sim A}|f_{n}-f|+\int_{A}|f_{n}|+\int_{A}|f|.
\end{align*}
令 $\varepsilon>0$. 根据 $\{f_{n}\}$ 的一致可积性, 存在 $\delta>0$, 使得对 $E$ 的任何满足 $m(A)<\delta$ 的可测子集 $A$, $\int_{A}|f_{n}|<\varepsilon/3$. 因此, 根据 Fatou 引理, 我们也有对 $E$ 的任何满足 $m(A)<\delta$ 的可测子集 $A$, $\int_{A}|f|<\varepsilon/3$. 由于 $f$ 是实值的而 $E$ 具有有限测度, Egoroff 定理告诉我们存在 $E$ 的可测子集 $E_{0}$, 使得 $m(E_{0})<\delta$ 且在 $E\sim E_{0}$ 上一致地 $\{f_{n}\}\to f$. 选取自然数 $N$, 使得在 $E\sim E_{0}$ 上对所有 $n\geqslant N$, $|f_{n}-f|<\varepsilon/[3\cdot m(E)]$. 在积分不等式中取 $A=E_{0}$. 若 $n\geqslant N$, 则
\begin{align*}
\left|\int_{E}f_{n}-\int_{E}f\right|\leqslant\int_{E\sim E_{0}}|f_{n}-f|+\int_{E_{0}}|f_{n}|+\int_{E_{0}}|f|<\varepsilon/[3\cdot m(E)]\cdot m(E\sim E_{0})+\varepsilon/3+\varepsilon/3\leqslant\varepsilon.
\end{align*}
这完成了证明.
\end{proof}

以下定理表明一致可积性的概念是确保下述合理性的实质性要素: 有限测度集上的非负函数序列 $\{h_{n}\}$ 在积分号下取极限逐点收敛到 $h\equiv 0$.

\begin{theorem}\label{theorem:ch04-一致可积判定}
令 $E$ 为有限测度集. 假定 $\{h_{n}\}$ 是在 $E$ 上 a.e. 逐点收敛于 $h\equiv 0$ 的非负可积函数序列. 则
\begin{align*}
\lim_{n\to\infty}\int_{E}h_{n}=0\quad\text{当且仅当}\{h_{n}\}\text{在 }E\text{ 上一致可积}.
\end{align*}
\end{theorem}

\begin{proof}
若 $\{h_{n}\}$ 是一致可积的, 则根据 Vitali 收敛定理, $\lim_{n\to\infty}\int_{E}h_{n}=0$. 反过来, 假定 $\lim_{n\to\infty}\int_{E}h_{n}=0$. 令 $\varepsilon>0$. 我们可以选取自然数 $N$, 使得当 $n\geqslant N$, $\int_{E}h_{n}<\varepsilon$. 因此, 由于在 $E$ 上每个 $h_{n}\geqslant 0$,
\begin{align}
\text{若 }A\subseteq E\text{ 可测且 }n\geqslant N,\text{则}\int_{A}h_{n}<\varepsilon.\label{align:ch04-尾部估计}
\end{align}
根据命题\cref{proposition:ch04-有限族一致}, 有限族 $\{h_{n}\}_{n=1}^{N-1}$ 在 $E$ 上是一致可积的. 令 $\delta$ 对应于关于 $\{h_{n}\}_{n=1}^{N-1}$ 一致可积性准则中的 $\varepsilon$ 挑战. 我们从\eqref{align:ch04-尾部估计}推出 $\delta$ 也对应于关于 $\{h_{n}\}_{n=1}^{\infty}$ 一致可积性准则的 $\varepsilon$ 挑战.
\end{proof}

\begin{exercise}\label{exercise:ch04-40}
令 $f$ 在 $\mathbf{R}$ 上可积. 证明对所有 $x\in\mathbf{R}$, 定义为
\begin{align*}
F(x)=\int_{-\infty}^{x}f
\end{align*}
的函数 $F$ 是正确定义的且是连续的. 它必定是 Lipschitz 的吗?
\end{exercise}

\begin{exercise}\label{exercise:ch04-41}
证明: 若 $E=\mathbf{R}$, 命题\cref{proposition:ch04-极限可积}不成立.
\end{exercise}

\begin{exercise}\label{exercise:ch04-42}
证明: 若没有 $h_{n}$ 是非负的假设, 定理\cref{theorem:ch04-一致可积判定}不成立.
\end{exercise}

\begin{exercise}\label{exercise:ch04-43}
令函数序列 $\{f_{n}\}$ 和 $\{g_{n}\}$ 在 $E$ 上一致可积. 证明对任何 $\alpha$ 和 $\beta$, 线性组合序列 $\{\alpha f_{n}+\beta g_{n}\}$ 在 $E$ 上也是一致可积的.
\end{exercise}

\begin{exercise}\label{exercise:ch04-44}
令 $f$ 在 $\mathbf{R}$ 上可积且 $\varepsilon>0$. 证明以下三个逼近性质.
\begin{enumerate}[(i)]
\item 存在 $\mathbf{R}$ 上的简单函数 $\eta$, 它具有有限支撑且 $\int_{\mathbf{R}}|f-\eta|<\varepsilon$. (提示: 首先对 $f$ 非负的情形证明.)
\item 存在 $\mathbf{R}$ 上的阶梯函数 $s$, 它在一个闭有界区间外消失且 $\int_{\mathbf{R}}|f-s|<\varepsilon$. (提示: 应用 (i) 部分与第 3 章的习题 18.)
\item 存在 $\mathbf{R}$ 上的连续函数 $g$, 它在一个有界集外消失且 $\int_{\mathbf{R}}|f-g|<\varepsilon$.
\end{enumerate}
\end{exercise}

\begin{exercise}\label{exercise:ch04-45}
令 $f$ 在 $E$ 上可积. 通过设在 $E$ 外 $\hat{f}\equiv 0$, 定义 $\hat{f}$ 为 $f$ 到整个 $\mathbf{R}$ 的延拓. 证明 $\hat{f}$ 在 $\mathbf{R}$ 上可积且 $\int_{E}f=\int_{\mathbf{R}}\hat{f}$.
用此与前一个习题的 (i) 与 (iii) 部分, 证明对 $\varepsilon>0$, 存在 $E$ 上的简单函数 $\eta$ 与 $E$ 上的连续函数 $g$, 使得
\begin{align*}
\int_{E}|f-\eta|<\varepsilon\quad\text{且}\quad\int_{E}|f-g|<\varepsilon.
\end{align*}
\end{exercise}

\begin{exercise}\label{exercise:ch04-46}
(Riemann-Lebesgue) 令 $f$ 在 $(-\infty,\infty)$ 上可积. 证明
\begin{align*}
\lim_{n\to\infty}\int_{-\infty}^{\infty}f(x)\cos nx\,dx=0.
\end{align*}
(提示: 首先对 $f$ 为闭有界区间外消失的阶梯函数证明, 接着用习题 44 的逼近性质 (ii).)
\end{exercise}

\begin{exercise}\label{exercise:ch04-47}
令 $f$ 在 $(-\infty,\infty)$ 上可积.
\begin{enumerate}[(i)]
\item 证明对每个 $t$,
\begin{align*}
\int_{-\infty}^{\infty}f(x)\,dx=\int_{-\infty}^{\infty}f(x+t)\,dx.
\end{align*}
\item 令 $g$ 为 $\mathbf{R}$ 上的有界可测函数. 证明
\begin{align*}
\lim_{t\to 0}\int_{-\infty}^{\infty}g(x)\cdot[f(x)-f(x+t)]=0.
\end{align*}
\end{enumerate}
(提示: 首先用 $f$ 在 $\mathbf{R}$ 上的一致连续性, 若 $f$ 是连续的且在一个有界集外消失的话, 证明等式. 接着用习题 44 的逼近性质 (iii).)
\end{exercise}

\begin{exercise}\label{exercise:ch04-48}
令 $f$ 在 $E$ 上可积而 $g$ 是 $E$ 上的有界可测函数. 证明 $f\cdot g$ 在 $E$ 上可积.
\end{exercise}

\begin{exercise}\label{exercise:ch04-49}
令 $f$ 在 $\mathbf{R}$ 上可积. 证明以下四个断言是等价的:
\begin{enumerate}[(i)]
\item 在 $\mathbf{R}$ 上, a.e. $f=0$.
\item 对 $\mathbf{R}$ 上的每个有界可测函数 $g$, $\int_{\mathbf{R}}fg=0$.
\item 对每个可测集 $A$, $\int_{A}f=0$.
\item 对每个开集 $\mathcal{O}$, $\int_{\mathcal{O}}f=0$.
\end{enumerate}
\end{exercise}

\begin{exercise}\label{exercise:ch04-50}
令 $\mathcal{F}$ 为函数族, 其中每个函数在 $E$ 上可积. 证明 $\mathcal{F}$ 在 $E$ 上一致可积, 当且仅当对每个 $\varepsilon>0$, 存在 $\delta>0$, 使得对每个 $f\in\mathcal{F}$,
\begin{align*}
\text{若 }A\subseteq E\text{ 是可测的且 }m(A)<\delta,\text{则}\left|\int_{A}f\right|<\varepsilon.
\end{align*}
\end{exercise}

\begin{exercise}\label{exercise:ch04-51}
令 $\mathcal{F}$ 为函数族, 其中每个函数在 $E$ 上可积. 证明 $\mathcal{F}$ 在 $E$ 上一致可积, 当且仅当对每个 $\varepsilon>0$, 存在 $\delta>0$, 使得对所有 $f\in\mathcal{F}$,
\begin{align*}
\text{若 }\mathcal{U}\text{ 是开的且 }m(E\cap\mathcal{U})<\delta,\text{则}\int_{E\cap\mathcal{U}}|f|<\varepsilon.
\end{align*}
\end{exercise}

\end{document}